\section{Data}
\label{sec:data}

The pipeline assembles eight public English datasets --- most published by the
Ministry of Housing, Communities and Local Government (MHCLG) or the Office for
National Statistics (ONS) --- into a single analysis-ready database, all joined on
the 2021-Census Lower-layer Super Output Area (LSOA; code \texttt{lsoa21},
33{,}755 in England). Table~\ref{tab:datasets} lists the sources. Every dataset is
published under the Open Government Licence (OGL) v3.0, with the single exception
noted below for address-level Energy Performance Certificate (EPC) records.

\begin{table}[htbp]
  \centering
  \caption{Public data sources for the pipeline. Grain is the native unit of
  observation; all are keyed or aggregated to the LSOA for modelling. Row counts
  are as ingested. MHCLG: Ministry of Housing, Communities and Local Government;
  ONS: Office for National Statistics; EPC: Energy Performance Certificate;
  IMD: Index of Multiple Deprivation; CQC: Care Quality Commission.}
  \label{tab:datasets}
  \small
  \begin{tabular}{@{}llrl@{}}
    \toprule
    Dataset & Grain & Rows & Role in model \\
    \midrule
    Dwelling fires (MHCLG) & incident & 451{,}172 & $P(\text{casualty}),\,P(S)$ \\
    Other-building fires (MHCLG) & incident & 231{,}426 & non-dwelling context \\
    Low-level geography (MHCLG) & incident & 8{,}502{,}911 & $P(I)$ spatial exposure \\
    EPC register (MHCLG) & certificate & 22{,}375{,}762 & building attributes ($x_b$) \\
    \quad --- latest per dwelling & dwelling & 17{,}302{,}650 & dwelling-stock aggregates \\
    IMD 2025 (MHCLG) & LSOA & 33{,}755 & deprivation covariates ($x_b$) \\
    Census 2021 (ONS, 5 tables) & LSOA & 33{,}755 & occupancy proxies ($x_b$) \\
    CQC active locations & premises & 56{,}870 & occupancy vulnerability \\
    Postcode $\to$ LSOA lookup (ONS) & postcode & 2{,}714{,}964 & EPC join key \\
    \bottomrule
  \end{tabular}
\end{table}

\paragraph{Sources.}
Fire incidence comes from the MHCLG fire statistics
collection \citep{mhclg2026firestats}: three record-level datasets and a suite of
aggregate FIRE-series benchmark tables. The \emph{dwelling-fires} dataset (one row
per primary dwelling fire attended by a fire and rescue service (FRS), FY
2010/11--2024/25) carries the ignition, alarm, spread and casualty fields used by
the consequence models; the \emph{low-level-geography} dataset locates every
incident to an LSOA and supplies the spatial exposure denominator. We ingested the
full 92-table FIRE catalogue and retain 61 benchmark tables, including FRS
response times (FIRE1001, used as a consequence covariate in
Section~\ref{sec:results-casualty}), building-type spread and casualty rates
(FIRE0301/0304, used to calibrate the institutional archetypes), and smoke-alarm
operation rates (FIRE0701/0703/0705). Building attributes come from the EPC
domestic register \citep{epcregister2025} --- 22.38\,M England lodgements,
deduplicated to 17.30\,M latest-per-dwelling certificates for stock aggregates
(built form, construction-age band, floor area, main fuel, tenure, floor level).
Local covariates come from the Index of Multiple Deprivation (IMD) 2025 (all seven
domains per LSOA) and five Census 2021 tables \citep{ons2021census}: accommodation
type (TS044), central heating (TS046), occupancy rating (TS052), tenure (TS054)
and household composition (TS003). Occupancy vulnerability is sharpened by the Care
Quality Commission (CQC) register of active health-and-social-care locations
(14{,}896 flagged as care homes) and, as a worked example of the
house-in-multiple-occupation (HMO) archetype, the London Borough of Camden HMO
licensing register (no national open HMO dataset exists).

\paragraph{Join architecture.}
The canonical spatial key is the LSOA. IMD, Census and the EPC dwelling-stock
aggregates join to it directly; MHCLG has re-coded the full incident history
(including 2010/11) to 2021 LSOAs, so no 2011$\to$2021 crosswalk is required. EPC
certificates carry a postcode but no LSOA, so they are joined via the ONS best-fit
UPRN join to the Ordnance Survey open identifier products, with postcode
best-fit as fallback (99.91\% of England lodgements matched, 98.6\% exactly
by UPRN). The record-level
dwelling-fire dataset carries no geography below the FRS area (withheld for
privacy), so those models are hierarchical over the 43--44 FRS areas rather than
over LSOAs; the spatial occurrence model instead uses the LSOA-resolved incident
counts against Census household denominators.

\paragraph{Validation.}
Twenty-seven automated checks pass (zero failures). Record-level dwelling-fire counts
reconcile to the published FIRE0201 aggregate at a ratio of 1.000 in every year
from 2010/11 to 2023/24 (residuals of a handful of incidents), and 0.995 in
2024/25 (a known Suffolk data gap, below). The 33{,}755-LSOA count matches exactly
across IMD and all five Census tables; 99.36\% of low-level-geography rows carry a
valid English LSOA that joins to IMD at 100\% coverage. As an independent
cross-check, the per-LSOA flat share derived from EPC agrees with the Census TS044
flat share at Pearson $r = 0.977$, confirming the postcode-to-LSOA join and the
EPC stock aggregation.

\paragraph{Known quirks (carried honestly into the models).}
Four data-quality issues are documented and handled rather than hidden.
(i)~\emph{Suffolk 2024/25 gap:} the dwelling dataset is missing Suffolk records for
September 2024--March 2025 (212 fires short), so 2024/25 record-level counts sit
just below the FIRE0201 aggregate and Suffolk is excluded from that year in the
occurrence model; the principled treatment --- marginalising the missing months
under a partial-observation likelihood rather than dropping the area --- is v2
work. (ii)~\emph{`Not known' LSOAs:} 54{,}194 incident rows carry an unresolved
LSOA code (a publisher over-count for three services) and are flagged and excluded
from spatial modelling via \texttt{is\_valid\_lsoa21}; treating the true LSOA as a
latent allocation to be inferred, rather than discarding these rows, is likewise
deferred to v2. (iii)~\emph{EPC \texttt{floor\_level} is flat-only:} it is
populated for flats/maisonettes (roughly 30\% of stock) and blank for houses, so
its 69\% null rate is structural, not missing data, and it is used only after
filtering to flats. (iv)~\emph{Address-level EPC provenance, by design:} the
processed database holds \emph{no} Royal Mail address strings --- only the Unique
Property Reference Number (UPRN), postcode, administrative geography and dwelling
attributes. Dwelling-level linkage is via the ONS UPRN Directory (ONSUD; exact
UPRN joins under OGL), with postcode best-fit as a fallback; the Royal Mail
intellectual property that carries the address-level licence restriction exists
only in the raw source, which we do not redistribute. Reproduction is therefore via
the documented pipeline plus a free EPC-service registration, and every EPC
quantity still enters the model as an LSOA-level aggregate.
